\documentclass[11pt]{article}
\usepackage[margin=1in]{geometry}
\usepackage{booktabs}
\usepackage{longtable}
\usepackage{amsmath}
\usepackage{hyperref}
\hypersetup{colorlinks=true, linkcolor=blue, urlcolor=blue}
\sloppy

\title{SF Planning Commission Minutes:\\ Pipeline Status and Preliminary Results}
\author{Dan Post}
\date{August 2026}

\begin{document}
\maketitle

\section{Purpose}

This directory documents the empirical backbone of the project: turning twenty-eight years
of San Francisco Planning Commission meeting minutes (1998--2026) into an \emph{item-level}
dataset of discretionary land-use decisions. One row is one request heard by the
Commission --- a conditional use, a discretionary review, a variance, a rezoning --- with
its case number, address, zoning district, staff recommendation, disposition, roll-call
vote, and hearing date.

That dataset is what makes the ``market for housing regulation'' measurable. It lets us
ask who obtains discretionary approval, how often staff recommendations are overridden,
how the approval rate moves with the political composition of the Commission, and how any
of this responds to state preemption. The theory side lives in
\texttt{output/political\_economic\_housing\_model/}; the extension of the same pipeline
to the rest of the Bay Area is documented in the recon directories.

\section{What is in this directory}

The directory is \textbf{one folder per memo}. Each holds \texttt{<name>.tex}, its compiled
\texttt{.pdf}, and its own \texttt{figures/} and \texttt{tables/}, so a memo and everything it
depends on move together; every figure and table is written by a script, and each memo's last
section names the script that regenerates it. \texttt{README.md} in this directory is the
index. The spec and reference documents the pipeline code links to by path live in
\texttt{notes/}.

\begin{longtable}{p{0.37\textwidth}p{0.55\textwidth}}
\toprule
Folder & What it answers \\
\midrule
\endhead
\texttt{memo/} & This memo: pipeline status and what is still open. Edited in place. \\
\texttt{meeting\_level\_info/} & How each item is matched to the meeting it was heard at, the
meeting-level fields, and their accuracy against the gold set. \\
\texttt{extraction\_method\_comparison/} & Regex against Claude, field by field, on a frozen
train/test split, with accuracy and over-extraction reported separately. \\
\texttt{discretionary\_review\_patterns/} & What the 16{,}199-item corpus contains:
composition of the docket, dispositions, delay, commissioners, geography, Planning Code
citations. \\
\texttt{permit\_linkage/} & How far the entitlement can be followed to a real building permit
in DBI's records, and what the match buys. \\
\texttt{conditions\_of\_approval/} & Whether the \emph{content} of the Commission's conditions
can be recovered, given that \texttt{conditions\_imposed} is only a yes/no flag. \\
\texttt{predicting\_delay/} & How much elapsed hearing time is forecastable from what a
developer can observe at the first hearing --- in sample and out. \\
\texttt{data\_acquisition/} & The external data the real-options frame needs: a parcel-year
risk set, filing dates and two clocks, the conditional-use-to-permit bridge, scale, prices and
fees. \\
\texttt{conditions\_content/} & A census of condition text for every case and motion number,
the earliest year with text by source, and what the Commission's conditions actually say. \\
\texttt{permits\_content/} & What DBI's permit record contains, and whether Commission-linked
projects differ from observably similar ones. \\
\texttt{datasf\_planning\_catalogue/} & Everything DataSF's catalogue returns for
``planning'', with join keys, staleness and judgement flags. \\
\texttt{notes/} & \texttt{labeling\_rules.md} (the coding manual),
\texttt{data\_infrastructure.md} (schema and flow), \texttt{minutes\_data\_availability.md},
\texttt{processing\_review.md}, \texttt{schema\_enrichment\_recommendation.md},
\texttt{hand\_label\_review\_guide.md}, \texttt{help\_string\_audit.md}. \\
\bottomrule
\end{longtable}

\section{The pipeline}

\textbf{Scrape} $\to$ \textbf{parse} $\to$ \textbf{date} $\to$ \textbf{label} $\to$
\textbf{extract}. Code lives in \texttt{code/commission\_minutes\_processing/}; the corpus
lives outside the repo on Dropbox and is reached through \texttt{paths.py}.

\emph{Parsing} splits each source document into blocks, one per agenda item, delimited by
case headers, numbered agenda items, section headers, adjournment markers, and
draft-minutes-adoption appendages. Two document eras are handled: scraped HTML
(1998--2014) and PDF (2015--present).

\emph{Dating} is a separate stage, and deliberately so. A block never states its own
hearing date --- the meeting header above it does --- so the date is a structural question
about which meeting a block falls under, not a textual one about the block. It is answered
by \texttt{assign\_meeting\_dates.py} and validated against hand-marked meeting boundaries
collected in \texttt{date\_boundary\_app/}. See the companion memo,
\texttt{meeting\_level\_info.tex}.

\emph{Meeting-level labelling} sits above the items. Time of day, whether the sitting was
regular, special, joint or closed, the room, who sat and who staffed it are properties of
the hearing, shared by every item heard at it, so they are recorded once per meeting and
joined on the date rather than repeated on 23{,}000 rows.
\texttt{meeting\_headers.py} cuts the header window around each known meeting start and
pre-fills twelve such fields for confirmation in the boundary app.

\emph{Labeling} is done in a local Flask app against the shared schema, with a machine
pre-fill the human corrects. \emph{Extraction} then either fine-tunes a model on those
labels or few-shot prompts one, and \texttt{learning\_curve.py} plots field accuracy
against label count so labeling stops at the empirical plateau rather than a guess.

\subsection{The schema is the single source of truth}

\texttt{extraction\_common.py} holds one list, \texttt{SCHEMA}, from which the labeling
form, the model prompt, the training builder's required keys, and the evaluation metric
are all generated. Editing that list propagates everywhere. It currently carries
\textbf{28 fields} in six groups: Identity, Location, Zoning \& scale, Process, Politics,
Conditions.

Four fields have been dropped as the schema settled, each for a stated reason:
\texttt{jurisdiction} (constant), \texttt{supervisorial\_district} (recoverable from the
address), \texttt{item} (agenda number, no analytic signal), and --- as of 30 August 2026
--- \texttt{meeting\_date}, which is a property of the meeting rather than the item and is
now supplied by the dating stage.

Two design decisions in the schema are worth stating because they carry the empirics.
First, \texttt{action} records the disposition \emph{family} (\texttt{approved},
\texttt{disapproved}, \texttt{continued}, \texttt{took\_dr\_and\_approved}, \ldots), with
``with conditions'' and ``as modified'' split out into two orthogonal flags. This removed a
three-way overlap in the old vocabulary and lets conditions attach to non-approval
outcomes. Second, the staff recommendation is stored twice: verbatim
(\texttt{preliminary\_recommendation}) and bucketed to the same disposition families
(\texttt{preliminary\_recommendation\_category}). The gap between that category and
\texttt{action} \emph{is} the staff-override signal the project studies, and storing both
makes ``was the staff recommendation followed?'' a direct enum comparison rather than
string wrangling.

\section{Preliminary results}

\subsection{Corpus}

\begin{table}[h]
\centering
\begin{tabular}{lr}
\toprule
Quantity & Count \\
\midrule
Parsed items, 1998--2026 & 25{,}927 \\
\quad HTML era (1998--2014) & 17{,}080 across 724 pages \\
\quad Modern era (2015--2026) & 8{,}847 \\
\quad of which carry a case number & 16{,}269 \\
Meetings identified, HTML era & 818 \\
Meetings identified, whole corpus & 1{,}219 across 1{,}125 documents \\
Pages holding more than one meeting & 41 (max 5) \\
\bottomrule
\end{tabular}
\end{table}

\subsection{Dating}

The dating stage re-dates every HTML-era item from its position in the source page.
Distinct hearing dates per year moved from 12--17 to 35--47 for 1998--2000, in line with the
Commission's roughly weekly sitting; later years were already close to correct. Against the
hand-marked gold standard --- now 116 documents, scored on HAND marks only, since a later
validation round pre-marked its documents by machine and those must not be scored as gold
--- boundary recall is \textbf{1.000} and precision \textbf{0.992}, positional recall
\textbf{1.000} and precision \textbf{0.984}, and \textbf{item date accuracy is 1.000 on
1{,}752 scored items}. All modern PDFs agree with their file names. Re-derive with
\texttt{python date\_boundary\_app/app.py --score}; details and caveats are in the companion
memo.

\subsection{Meetings}

\textbf{134 meetings} are hand-confirmed, and the extraction now runs over the whole
corpus rather than the sample: \textbf{1{,}219 meetings, 1998--2026}, written to
\texttt{meetings\_all.csv}. Against the hand-confirmed set the ten meeting-level fields
agree on \textbf{93.9\%} of 938 values (\texttt{extract\_all\_meetings.py --score}); the
out-of-sample generalisation estimate from the last frozen round is \textbf{96.0\%}. The two
differ because the first covers every round under a stricter comparison and is
unadjudicated; the meeting-level memo keeps them apart. The sample is already informative about the panel's
composition: 90 regular sittings, 7 special, and 5 joint --- with the Redevelopment Agency
Commission (1998), Building Inspection (2007, 2018), Public Health (2011) and Historic
Preservation (2023). Two joint sittings are headed simply ``Special Meeting'' and are
identifiable as joint only from the two bodies named above the date.

\subsection{Labeling}

\begin{table}[h]
\centering
\begin{tabular}{lrl}
\toprule
Status & Items & Meaning \\
\midrule
\texttt{done} & 289 & confirmed gold \\
\texttt{flagged} & 89 & a specific issue to recheck \\
\texttt{not\_an\_item} & 9{,}658 & agenda scaffolding, not a land-use decision \\
\texttt{todo} & 15{,}891 & unlabeled \\
\bottomrule
\end{tabular}
\end{table}

The queue is ordered rare-class-first and interleaved by year, so labeling top-down fills
scarce request types and dispositions evenly across the panel rather than exhausting 1998
first. A fixed random sample per year is flagged separately so the gold set is
representative of the whole period rather than of whichever years were labeled earliest.

\section{Status and open questions}

\begin{itemize}
\item \textbf{Settled.} The schema (v2); the parser's block boundaries; the dating stage
and its validation; the labeling workflow and coding manual; the extraction configuration
(Haiku 4.5, six era-matched examples, structured outputs with evidence spans) and its
corpus-wide run.
\item \textbf{In progress.} Item-level hand-labeling; 289 items are confirmed gold and the
frozen split is 155 train / 139 test. The unified review queue holds the remaining work.
The meeting-level extraction is complete over the whole corpus (1{,}219 meetings) with 134
hand-confirmed.
\item \textbf{Open.} The adjudication queue still holds 143 unjudged gold-vs-model
disagreements, and they are built against a superseded zero-shot run rather than the live
configuration --- so the share of measured ``model error'' that is really gold error is
known only for the batch already judged (54\% model, 43\% gold, 3\% both). The learning
curve has not been run at full size, so the number of labels required for a local model is
still an estimate. Six commissioners' roll-call entries still merge two surnames where the
minutes omit a separator (`Melgar Moore', `Moore Sugaya'); splitting them needs a roster
the project does not yet hold.
\item \textbf{Decision pending.} Whether to extend meeting-level hand-confirmation beyond
the 134 already done. The pre-fill is good enough that the cost is confirmation rather than
transcription, but the analytic payoff of the meeting-level fields is not yet established.
\item \textbf{Known non-transfer.} The heuristic extractor is tuned to SF minutes and does
not transfer to other jurisdictions; extending the pipeline requires re-testing extraction
on real documents with human review, not unattended.
\end{itemize}

\end{document}
